\section{Finite coverage by long experiments}
\label{sec:v77-coverage}

The conclusion of Theorem~\ref{thm:v77-main} becomes a whole-table
statement only when the experiment genuinely covers the boundary.
We prove the existence of such experiments on an explicit class of
periodic smooth dispersing tables.  An upper bound on a collision-free
flight will be called a finite-horizon bound.  Obstacle separation and
positive curvature are strict throughout this section.

\begin{lemma}[Regular extension of a fan]\label{lem:v77-fan}
Fix a boundary point $x$ in a finite-horizon dispersing table and an
open interval of nongrazing directions entering the free region.
For every integer $J$, this interval contains a direction whose forward
trajectory is regular for its first $J$ collisions.  The same statement
holds for backward trajectories.  In particular, a clear nongrazing
one-flight patch can be restricted about a point on either endpoint
boundary so that it has a regular prefix of any prescribed finite
length, together with positive adjacent free-time margins.
\end{lemma}
\begin{proof}
At the first flight, only finitely many obstacle lifts are within the
finite-horizon bound.  From a fixed exterior point a strictly convex
body has at most two tangent rays.  Removing these directions leaves
open regular branches.  On each such branch the outgoing family is a
diverging ray fan.  Here is the transversality fact needed to iterate
this argument.  Write $Y$ for transverse separation of neighboring rays
and $\mathcal B$ for wave-front curvature.  Free propagation over length
$\tau$ and a dispersing reflection give respectively
\begin{equation}\label{eq:v77-fan-curvature}
 \mathcal B_{\rm free}=\frac{\mathcal B}{1+\tau\mathcal B},
 \qquad
 \mathcal B_{\rm reflected}=\mathcal B_{\rm incoming}
                                  +\frac{2\kappa}{\nu}.
\end{equation}
These formulas follow by differentiating the direction of the normal
rays to a wave front and then the specular reflection identity.  A point
fan has $Y=\tau$ at its first positive distance, in the initial-angle
normalization; thereafter free propagation multiplies a positive
separation by $1+\tau\mathcal B$, and reflection preserves its
nonvanishing, with the corresponding incidence-coordinate factor.
Thus $Y$ never vanishes before a singular collision.  This is also the
no-conjugacy calculation encoded by the positive form
\eqref{eq:v77-hessian}.

At a tangent encounter with a possible next obstacle, let $F$ be a
smooth local defining function for that obstacle and let
$r(\theta,t)$ be the ray family.  At tangency, $F_t=0$ and
$F_{tt}\ne0$ by positive curvature.  Solving $F_t=0$ for the
location of closest approach, the derivative of its signed distance
with respect to the fan parameter is a nonzero multiple of $Y$:
the normal to the obstacle is perpendicular to the tangent ray.
Hence a tangent direction is an isolated zero on every previously
regular parameter branch.  There are countably many such branches,
and finitely many next obstacle lifts on each bounded flight.  At each
fixed collision order only a countable set of additional directions is
excluded.  A finite union of these sets cannot exhaust an open interval.
Finite horizon prevents a regular ray from escaping without a next
collision.  This proves the forward assertion; reversal proves the
backward assertion.

For a given visible patch and fixed point $x$ at its initial end,
directions reaching its other open endpoint arc contain an open interval.
At $x$, reflection converts them smoothly to backward directions.
Choose one avoiding the first $J+1$ backward singularities.  The first
$J$ flights supply a prefix and the extra backward flight supplies the
initial free-time margin.  The terminal outgoing free segment has a
positive length.  Restricting the regular trajectory to compact small
coordinate boxes preserves all these margins and keeps $x$ inside the
final-flight patch.
\end{proof}

\begin{lemma}[Connected visibility graph of the lifted obstacles]
\label{lem:v77-visible-connected}
The graph whose vertices are all obstacle lifts and whose edges are
clear nongrazing one-flight patches is connected.
\end{lemma}
\begin{proof}
Suppose its vertices were partitioned into a connected component and
its nonempty complement.  The minimum distance between a body in one
part and a body in the other is attained.  Indeed, start with any
cross-part pair at distance $D$; modulo the finite list of obstacle
orbits there are only finitely many relative lattice translations whose
body distance is at most $D$.  Their distances form a finite nonempty
list, even if the partition is not periodic.

Let $[x,y]$ realize the least cross-part distance.  If another body
met its open segment, that body would be in one of the two parts and
would be strictly closer to the endpoint body in the other part,
a contradiction.  The open segment is therefore clear.  At the
minimizing endpoints it is normal to both smooth strictly convex
boundaries, so it is nongrazing.  Compactness and local finiteness give
clearance from all other obstacles, and small endpoint perturbations
produce a one-flight patch.  This is an edge between the two parts,
contradicting the choice of a graph component.
\end{proof}

\begin{theorem}[Finite identifying experiments and their persistence]
\label{thm:v77-finite-atlas}
Let a finite-horizon periodic dispersing table have locally nonconic
$C^\infty$ boundaries.  For every integer $J\ge1$ it admits a finite
identifying atlas for its complete periodic table, in which every
measured prefix has at least $J$ flights and every measured extension
one additional flight.  Its scalar coordinate domains, branch and lift
labels, anchor parameters and three absolute clocks per branch can be
fixed on a nonempty finite-smooth open neighborhood of the table.
Every realization of the same admissible atlas data is congruent to
it, whether or not that realization is close to it.
\end{theorem}
\begin{proof}
Choose one lift of each quotient obstacle and, in addition, the two
basis translates of one of them.  Lemma~\ref{lem:v77-visible-connected}
provides finitely many graph paths connecting this finite set.  Add the
finitely many intermediate lifts on these paths to the set of target
bodies.  Choose one clear patch on each path edge.  Slight perturbation
within a patch makes the backward prefix regular for the required
length by Lemma~\ref{lem:v77-fan}, without changing its two obstacle
labels or destroying its clearance.

For each point of each target boundary choose a nongrazing outgoing
ray with a regular next collision and a regular backward prefix of
length at least $J$.  Such rays exist by Lemma~\ref{lem:v77-fan}.
Lemmas~\ref{lem:v77-regular-action} and
Proposition~\ref{prop:v77-cancellation} give an open final-flight
rectangle whose first boundary interval contains that point, together
with two branch gates and a unique scalar matching bracket.  The
initial final-flight interval is nonconic by hypothesis.  Lemma
\ref{lem:v77-nonconic-certificate} supplies six and three interior
anchors for that rectangle.  Compactness of the finitely many target
boundaries selects finitely many such intervals covering them.  Choose
slightly smaller open covering intervals with closures in the constructed
rectangles, preserving the anchors in the larger patches as necessary.

A finite open cover of a connected closed boundary has a connected
overlap graph after redundant disconnected components are removed:
otherwise the unions of its graph components would separate the
boundary into disjoint nonempty open sets.  In each nonempty open
overlap choose three distinct parameter values.  The independently
reconstructed patches therefore register all pieces of each target
boundary.  The previously selected path-edge patches connect these
registrations across all target bodies.  Their second endpoint patches
and any further auxiliary patches attach to an already covered target
piece and do not disconnect the graph.  On the two translated copies
include a common scalar subinterval from the atlas transported by the
deck labels.  The resulting registration graph and target satisfy
Definition~\ref{def:v77-atlas}.

There are finitely many branch gates.  Shrink them, when needed, before
choosing the finite subcover and anchors, and apply Lemma
\ref{lem:v77-clock-design} to choose clock triples with strict margins.
Nonconstancy of each action supplies a usable pair of scalar anchors
in Theorem~\ref{thm:v77-clock}.  All choices are finite.  In particular
no limiting sequence of patch radii is being used to turn density of
visited points into coverage.

The positive Hessians, nonzero twists, clock and root brackets,
noncollinear overlap triples, and rank certificates all have positive
margins on these finitely many compact domains.  The finite implicit
construction and continuity preserve them on a $C^r$ neighborhood,
with $r$ chosen large enough for the finitely stipulated derivatives.
The scalar covering intervals and overlap maps are fixed and still
cover the perturbed boundary parametrizations.  Thus the same design
is an identifying atlas on this open class.  The last assertion is
Theorem~\ref{thm:v77-main}, whose uniqueness argument uses the
certificates and not nearness of the candidates.
\end{proof}

This theorem establishes existence and persistence of a finite design.
Its constants and the number of settings may depend on the table and
on $J$.  It does not replace the known scalar boundary atlas or exact
branch labels by conclusions.  Nor does it assert that finitely many
sampled points determine an arbitrary smooth boundary: each clock
setting here specifies a whole conditional law.

\subsection{An explicit open family with nonanalytic members}

\begin{proposition}[Nonempty smooth whole-table classes]
\label{prop:v77-example}
There are finite-horizon periodic tables to which Theorem
\ref{thm:v77-finite-atlas} applies, and their persistent identifying
atlas classes contain nonanalytic smooth tables.
\end{proposition}
\begin{proof}
Use the triangular lattice generated by $(1,0)$ and
$(1/2,\sqrt3/2)$.  Choose $R,\varepsilon$ with
\[
 \frac{\sqrt3}{4}<R-|\varepsilon|,
 \qquad R+|\varepsilon|<\frac12,
 \qquad 8|\varepsilon|<R,
 \qquad \varepsilon\ne0.
\]
At each lattice point place the strictly convex body with support
function
\begin{equation}\label{eq:v77-example-support}
 h(\theta)=R+\varepsilon\cos(3\theta).
\end{equation}
Its radius of curvature is $h+h''=R-8\varepsilon\cos(3\theta)>0$.
It contains the centered disk of radius $R-|\varepsilon|$ and is
contained in the disk of radius $R+|\varepsilon|$, so its translates
are disjoint.

We verify finite horizon for the contained disks.  An irrational
direction on the torus gives a dense straight line and hence meets a
disk interior.  In a rational direction parallel to a primitive lattice
vector $v$, the rows of lattice centers have spacing
$\sqrt3/(2|v|)\le\sqrt3/2$.  Every parallel line is within
$\sqrt3/4$ of one row and meets a disk interior since the disk radius
is larger.  Thus no full line avoids the disk interiors.  If there were
arbitrarily long free segments, translate their midpoints into a fixed
cell and pass to a subsequence of points and directions.  The limit
would be a complete line avoiding all disk interiors, a contradiction.
The larger bodies therefore have finite horizon as well.

The boundary in \eqref{eq:v77-example-support} is analytic, has
threefold rotational symmetry and is not a circle.  It cannot lie in
a conic: a conic containing a closed strictly convex curve must be an
ellipse, and an ellipse invariant under a rotation by $2\pi/3$ is a
circle.  If any open boundary arc lay in a conic, analyticity of the
parametrization would make the corresponding quadratic polynomial
vanish on the whole boundary.  Thus the boundary is locally nonconic.
Theorem~\ref{thm:v77-finite-atlas} applies.

Fix one of its finite identifying atlases.  Modify the support function
by a sufficiently small $C^\infty$ compactly supported bump which is
flat but not identically zero at an endpoint of its support.  Its size
can be arbitrarily small in any prescribed finite $C^r$ norm.  Positivity
of $h+h''$, disk containment, separation, and all finite atlas margins
persist.  The new support function is not analytic, while the same
finite experiment still determines its complete boundary by actual
function inversion.  No analytic continuation is applied to this
perturbed table.
\end{proof}
